% Nouveau EGA V (cote 25): archive cote 25, site page 94 (PDF page 95).
% typescript. Draft compared: none.
% AI-transcribed by Claude Opus 5.5 (claude-opus-5-5, adaptive thinking, effort high), Message Batches API, batch ega25b (msgbatch_01A7n5YuXubUPNK1zsepi8RG), 2026-11-13; not yet checked by eye.
\scanpage{25/94}
\mspage{88}
semble qu'elle doive faire un axiome primitif de la situation (dans la
situation du n°2.2. , elle \unsure{résultait} d'un argument global assez \unsure{spécifil|spécifiq} ):
\uline{pour toute partie fermée Zrréductible} Z \uline{de} P d\uline{e dimension} m, $\underline{Z}_{\hand{\eta}} \neq \emptyset$ \var{\xout{xxx}}{}{typewriter strike-out}{ega25_25_094-v0001}
\var{\xout{pour tout La précision sur la dimension sera alors conséquence facile}}{}{typewriter strike-out}{ega25_25_094-v0002}
\var{\xout{de l'hypothèse suivante, de nature différentielle: P est lisse sur S,}}{}{typewriter strike-out}{ega25_25_094-v0003}
\var{\xout{de dimension $\underline{P}$ $P \times_S T$ une immersion, et pour tout}}{}{typewriter strike-out}{ega25_25_094-v0004} Reprenant alors la
partie fermée Z de X,\var{\xout{de di}}{}{typewriter strike-out}{ega25_25_094-v0005} telle que dim f(Z) $\hand{\geqslant}$ m, on voit que $\underline{Z}_{\hand{\eta}} \rightarrow G$ est
dominant, et par suite $\underline{Z}_{\hand{\eta}}$ est de dimension dim $\underline{Z}$ - dim G = dim Z-m.

Ces propriétés permettent de développer dans le contexte présent
les résultats correspondant à 2.1. à 2.11. Il faut une condition sur (D)
assurant la validité de 2.12. i.e. que si X est lisse, alors $\underline{X}$ l'est
\unsure{éggalement|également}, en supposant f:X $\longrightarrow$ P non ramifié. N$_{\mathrm{o}}$us supposerons
\unsure{maàntenant|maintenant} P \uline{lisse sur} k,\var{\inserted{$\underline{P} \rightarrow P \times_k G$ quasi-fini,}}{$\underline{P} \rightarrow P \times_k G$ quasi-fini,}{insertion}{ega25_25_094-v0006} et la condition suivante vérifiée\var{\inserted[typed]{(où nous supposons k alg clos)}}{(où nous supposons k alg clos)}{insertion}{ega25_25_094-v0007}: pour
tout $x \in P(k)$, et tout sous-espace vectoriel V \var{\inserted[typed]{de dimension $n \geqslant m$}}{de dimension $n \geqslant m$}{insertion}{ega25_25_094-v0008} de l'espace tangent
$T_x$(P) de P en x, considérons l'ensemble\var{\inserted{E(x,V)}}{E(x,V)}{insertion}{ega25_25_094-v0009} \var{\xout{Ex}}{}{typewriter strike-out}{ega25_25_094-v0010} des $\hand{\xi} \in G(k)$ tels que
$\underline{P}_{\hand{\xi}}$ ait un point\var{\inserted{\unsure{z}}}{\unsure{z}}{insertion}{ega25_25_094-v0011} au dessus de x \var{\xout{en lequel on ait bien $\underline{P}$ P soit ramifié}}{}{typewriter strike-out}{ega25_25_094-v0012}
ne satisfaisant \uline{pas} l'ensemble de conditions suivantes: $\underline{P}_{\hand{\xi}}$ est lisse en
z, \var{\xout{le morphisme $\underline{P}$ P est non ramifié en z, et l'image de l'espace}}{}{typewriter strike-out}{ega25_25_094-v0013}
l'application tangente à $\underline{P}_{\hand{\xi}} \rightarrow P$ en z : $T_z(\underline{P}\ ) \rightarrow T_x(P)$ est injective
(i.e. $\underline{P}_{\hand{\xi}} \rightarrow P$ non ramifié en z) et son image est ``transversale'' à V
i.e. sa somme avec V est $T_x(P)$. \uline{Alors} E(x,V) (dont nous savons que
c'est la trace d'un ensemble \var{\xout{constructible}}{}{typewriter strike-out}{ega25_25_094-v0014} \var{\inserted[typed]{\var{\deleted{fermé}}{}{deletion}{ega25_25_094-v0016}}}{\var{\deleted{fermé}}{}{deletion}{ega25_25_094-v0016}}{insertion}{ega25_25_094-v0015}\var{\inserted{\unsure{constructible}}}{\unsure{constructible}}{insertion}{ega25_25_094-v0017} bien déterminé de G sur G(k))
\uline{est de dimension} $\hand{\leqslant}$ N-n\unsure{\var{\over{+}{-}}{-}{writing over}{ega25_25_094-v0018}}1 . - Moyennant cette condition, \var{\xout{xx}}{}{typewriter strike-out}{ega25_25_094-v0019} l'appli\unsure{:}
cation du critère jacobien et un comptage de dimensions montre que
l'ensemble\var{\inserted{fermé E}}{fermé E}{insertion}{ega25_25_094-v0020} des poihts\var{\inserted{x}}{x}{insertion}{ega25_25_094-v0021} de $\underline{X}$ tels que $\underline{X} \longrightarrow G$ soit non lisse en x, \uline{ou}
$\underline{P} \longrightarrow$ \unsure{\var{\over{G}{G}}{G}{writing over}{ega25_25_094-v0022}} non lisse en $\underline{f}(x)$, ou $\underline{P} \longrightarrow P$ ramifié en $\underline{f}(x)$, est de
dimension $\hand{\leqslant}$ n + (N-n\unsure{\var{\over{+}{-}}{-}{writing over}{ega25_25_094-v0023}}1) = N-1 , \hand{(}X étant lisse partout de dimension n\hand{)}.

\var{\margin{94}}{\placedmargin{94}}{marginal note placed}{ega25_25_094-v0024}
